%% ma: L10-B14-0043
%% lop: 10
%% chuong: 5
%% bai: 14
%% dang: 10.14.8
%% loai: TL
%% muc: TH
%% dap_an: "-"
%% nguon: "chuyen-de-cac-so-dac-trung-cua-mau-so-lieu-khong-ghep-nhom-toan-10-chuong-trinh-moi.pdf (GV. Phan Nhật Linh) (P05), Bài 03 Dạng 2 Bài tập tự luận 3"
%% kien_thuc: "Lớp 10 Bài 13 (số trung bình), Bài 14 (phương sai, độ lệch chuẩn)"
%% trang_thai: chinh-thuc
%% ghi_chu: "Tự duyệt theo yêu cầu giáo viên 22/10/2026."
\begin{ex}
Sản lượng lúa (tạ) của $40$ thửa ruộng thí nghiệm có cùng diện tích được trình bày trong bảng tần số sau:
\begin{center}
\begin{tabular}{|c|c|c|c|c|c|}\hline
Sản lượng & $20$ & $21$ & $22$ & $23$ & $24$\\\hline
Tần số & $5$ & $8$ & $11$ & $10$ & $6$\\\hline
\end{tabular}
\end{center}
Tính phương sai của mẫu số liệu.
\loigiai{$\overline x=\dfrac{5\cdot20+8\cdot21+11\cdot22+10\cdot23+6\cdot24}{40}=22{,}1$.

$s^2=\dfrac{5(-2{,}1)^2+8(-1{,}1)^2+11(-0{,}1)^2+10(0{,}9)^2+6(1{,}9)^2}{40}=\dfrac{61{,}6}{40}=1{,}54$.}
\end{ex}
